% Owner: B
We ask which communities and connecting nodes emerge from graphs built
on the ingredients of sampled recipes and countries. We project the CulinaryDB
recipe--ingredient bipartite graph~\cite{culinarydb} onto ingredients and
weight retained edges by pointwise mutual information (PMI). We also
build a country graph from a dataset of national dishes~\cite{wunderlich2022data},
using Jaccard similarity between ingredient sets as edge weights. We use
graph statistics, degree, strength, betweenness, Louvain communities and
modularity. We examine articulation points and bridges, compare rankings
with Spearman correlation, and add a shuffle test for agreement with UN
regions. We examine how thresholds change connectivity and modularity.
The ingredient graph has five communities (five in 28 of 30 Louvain
seeds) and modularity $0.400$. Hop betweenness follows frequency ($\rho=0.968$), with Water,
Salt and Butter leading. With $1/\mathrm{PMI}$ distances, Couscous,
Lasagna and Shallot lead, and the correlation is $\rho=0.323$; a second
conversion preserves several leaders but changes their order. Country
communities agree with UN regions for $64.3\%$ of countries, against a
mean of $34.0\%$ over $1{,}000$ shuffles ($p\approx0.001$). We find that
the distance choice changes ingredient broker rankings, while country
communities show a partial association with geography.
